\chapter{Módulos e \texttt{source}}

\section{O que é um módulo em Hayashi}

Em \hayashi{}, um módulo é simplesmente um arquivo \texttt{.hay}. Não há sistema de
pacotes nem declarações de exportação — qualquer definição no arquivo fica disponível
no escopo de quem o carrega.

Este modelo é intencional: a linguagem está em crescimento, e a complexidade de um
sistema de módulos com namespaces, visibilidade e resolução de dependências é adicionada
quando há demanda real. Por ora, \lstinline{source} cobre a totalidade dos casos de
uso práticos.

\section{O comando \texttt{source}}

\lstinline{source} executa um arquivo \texttt{.hay} no contexto atual. Todas as
definições do arquivo ficam visíveis a partir desse ponto:

\begin{lstlisting}[caption={Carregando um módulo}]
source "lib/estatistica.hay"
source "lib/graficos.hay"

// funções de estatistica.hay e graficos.hay estão disponíveis
let resultado = calcular_gini(df)
\end{lstlisting}

\section{Organizando um projeto}

A estrutura típica de um projeto \hayashi{}:

\begin{lstlisting}[style=inline]
projeto/
  main.hay           -- ponto de entrada
  lib/
    dados.hay        -- funções de limpeza e carga
    estimadores.hay  -- funções econométricas customizadas
    tabelas.hay      -- formatação de resultados
  dados/
    raw.csv
    painel.dta
  resultados/
\end{lstlisting}

\texttt{main.hay}:

\begin{lstlisting}
source "lib/dados.hay"
source "lib/estimadores.hay"
source "lib/tabelas.hay"

load "dados/painel.dta" as df

let df_limpo = preparar(df)
let modelo   = estimar_fe(df_limpo, renda ~ educ + exp)

tabela_regressao(modelo, "resultados/tabela1.tex")
\end{lstlisting}

\section{Escopo de \texttt{source}}

O arquivo carregado com \lstinline{source} executa no escopo global do arquivo
principal. Isso significa:

\begin{itemize}
  \item Funções definidas no arquivo são globalmente visíveis.
  \item Variáveis definidas no nível superior do arquivo também ficam no escopo global.
  \item Não há namespaces: se dois arquivos definem \lstinline{fn calcular(...)}, o
        segundo sobrescreve o primeiro.
\end{itemize}

\begin{aviso}
  Por não haver namespaces, evite nomes genéricos em módulos reutilizáveis.
  Prefira \texttt{gini\_calcular} a apenas \texttt{calcular}.
\end{aviso}

\section{Criando uma biblioteca}

Um arquivo de biblioteca é um arquivo \texttt{.hay} com funções e constantes, sem
código de nível superior que produza efeitos colaterais na carga:

\begin{lstlisting}[caption={lib/utils.hay}]
const PI = 3.14159265358979

fn arredondar(x, casas) {
  let fator = 10^casas
  int(x * fator + 0.5) / float(fator)
}

fn entre(x, lo, hi) {
  x >= lo and x <= hi
}

fn clamp(x, lo, hi) {
  if x < lo { lo } else if x > hi { hi } else { x }
}
\end{lstlisting}

\begin{lstlisting}[caption={main.hay}]
source "lib/utils.hay"

display arredondar(3.14159, 2)    // 3.14
display entre(5, 1, 10)          // true
display clamp(15, 0, 10)         // 10
\end{lstlisting}

\section{Roadmap}

Em versões futuras, \hayashi{} terá um sistema de módulos com:

\begin{itemize}
  \item Namespaces explícitos (\texttt{use modulo::funcao})
  \item Distribuição via registro central
  \item Declarações de dependência em manifesto
\end{itemize}

A transição de \lstinline{source} para esse sistema será sem quebra de
compatibilidade: arquivos \texttt{.hay} continuarão funcionando como módulos.
